\chapter{ورود و شروع تحصیل}
\label{ch:arrival}

\section{روز ورود — فرودگاه}

\begin{enumerate}
  \item مدارک را در دسترس (نه در چمدان تحویل‌شده) نگه دارید
  \item فرم‌های ورود (در صورت نیاز) را با دقت پر کنید
  \item به سؤالات مرزبان کوتاه و شفاف پاسخ دهید
  \item آدرس محل اقامت و تماس دانشگاه را بدانید
\end{enumerate}

\begin{warningbox}
  اگر \lr{LOA}، \lr{CAS} یا \lr{I-20} شما مشروط است،
  مدارک تکمیل‌شده (مثلاً مدرک زبان نهایی) را همراه داشته باشید.
\end{warningbox}

\section{هفته اول}

\begin{checklistbox}
  \begin{checklist}
    \item ثبت‌نام در \lr{Orientation} دانشگاه
    \item دریافت کارت دانشجویی
    \item افتتاح حساب بانکی
    \item خرید سیم‌کارت محلی
    \item ثبت آدرس در اداره مهاجرت (در صورت نیاز)
    \item آشنایی با حمل‌ونقل عمومی
  \end{checklist}
\end{checklistbox}

\section{شناسه‌های مهم (بسته به کشور)}

\begin{table}[htbp]
  \centering
  \caption{شناسه‌های رایج}
  \label{tab:ids}
  \begin{tabular}{@{}lll@{}}
    \toprule
    \textbf{کشور} & \textbf{شناسه} & \textbf{کاربرد} \\
    \midrule
    کانادا & \lr{SIN} & کار دانشجویی، مالیات \\
    آمریکا & \lr{SSN} & کار (محدود برای \lr{F-1}) \\
    بریتانیا & \lr{NI Number} & کار پاره‌وقت \\
    \bottomrule
  \end{tabular}
\end{table}

\section{کار دانشجویی}

\begin{tipbox}
  ساعات مجاز کار در ویزای تحصیلی محدود است.
  قوانین را از سایت رسمی مهاجرت همان کشور بخوانید —
  کار بیش از حد مجاز می‌تواند ویزا را به خطر بیندازد.
\end{tipbox}

\section{سلامت روان و جامعه}

\begin{itemize}
  \item به مرکز مشاوره دانشگاه (\lr{Counselling}) سر بزنید
  \item در انجمن‌های دانشجویی ایرانی/بین‌المللی عضو شوید
  \item احساس تنهایی در ماه‌های اول طبیعی است — از کمک نترسید
\end{itemize}
